\section{The arithmetic foundation}\label{sec:arith}

\subsection{The exact version of iPRA}
We use first-order intuitionistic logic with equality, a symbol for every
primitive-recursive definition, its defining equations, successor
injectivity, $Sx\ne0$, and induction for quantifier-free formulas with
parameters. This is the system explicitly defined by
\texttt{PR}, \texttt{PTm}, and \texttt{PRAAx} in ProveIt's
\texttt{ConciseIPRA/PR.lean} \cite{repo-pr}.
Quantifier-free induction is not induction for arbitrary formulas about
geometric objects or real names.

Natural-number equality is decidable in this theory. Define its Boolean
test by primitive recursion, first using the zero test and predecessor,
and then comparing successors. The defining equations and
quantifier-free induction prove the test correct. Consequently every
finite Boolean combination of the integer equalities below permits
intuitionistically valid case analysis. This uses a decision procedure
for the specified predicate, rather than excluded middle for all formulas.

\subsection{Signed integers without a quotient-set assumption}
An integer is represented by a pair $(p,n)$ of natural numbers, read as
$p-n$. There is no need to form a set-theoretic quotient internally. Set
\begin{align*}
 (p,n)\simeq(p',n')&\quad\Longleftrightarrow\quad p+n'=p'+n,\\
 (p,n)+(r,s)&=(p+r,n+s),\\
 -(p,n)&=(n,p),\\
 (p,n)(r,s)&=(pr+ns,ps+nr).
\end{align*}
Zero and one are $(0,0)$ and $(1,0)$. These operations and the equality
test are primitive recursive and respect $\simeq$. Subsequent integer
equalities abbreviate this numerical equality relation. A tuple may be
left as finitely many variables or packed by a primitive-recursive
pairing function.

For completeness, the elementary arithmetic behind this interpretation
is as follows. Define $x+0=x$, $x+S y=S(x+y)$, $x\cdot0=0$, and
$x\cdot S y=x\cdot y+x$. Quantifier-free induction proves
$0+x=x$, $Sx+y=S(x+y)$, and associativity of addition. The first two
equations give commutativity of addition by induction on its second
argument. Induction on the repeated summand proves distributivity;
using distributivity, induction proves associativity and commutativity
of multiplication. Each induction formula is an equation with the other
quantities as parameters. Cancellation of a common additive summand
follows by induction on that summand and successor injectivity.
Expanding the pair operations now proves all commutative-ring laws:
both sides reduce to equal sums of natural monomials. The same expansions
show that replacing an input by an equivalent pair preserves the output.

We also need that these interpreted integers have no zero divisors.
Truncated subtraction $x\mathbin{\dot-}y$ is primitive recursive; its
recursion on $y$ uses predecessor. Its elementary defining identities
show that $(p,n)$ is equivalent to
\[
 (p\mathbin{\dot-}n,\ n\mathbin{\dot-}p),
\]
and that at least one of these two components is zero. One direct proof
uses the primitive-recursive comparison: strip successors from both
inputs until a zero is reached; the process can be bounded by $p+n$.
Induction on this explicit counter, with equations for the two registers
and the Boolean comparison result, is quantifier free. It proves the
usual alternatives $p=n+d$ or $n=p+d$, where the appropriate $d$ is the
displayed truncated difference. In either alternative the normalization
and its equality to $(p,n)$ follow by additive cancellation.

Every nonzero integer therefore has either the form $(S r,0)$ or
$(0,S r)$. The product of two such integers has a nonzero component,
since
\[
 S r\cdot S s=S\bigl(r+s+rs\bigr).
\]
This is unequal to zero by successor separation. Splitting on the
decidable zero test thus proves
\begin{equation}\label{eq:cancel}
 uv=0\ \land\ u\ne0\quad\Longrightarrow\quad v=0.
\end{equation}
All normalizations and comparisons here have primitive-recursive
definitions. None searches without a bound, and all the indicated
induction formulas are quantifier free.

\subsection{Why integer triples cover rational geometry}
A rational number can be presented by a signed numerator and a positive
natural denominator. Multiplying a triple by the product of its three
denominators produces an integer triple. Applying the same operation
to the six conic coefficients produces an integer quadratic equation.
No greatest common divisor or canonical representative is needed.
Multiplication by these nonzero denominators preserves incidence and
the zero or nonzero status of every determinant and cross product in
the statements. These operations are primitive recursive.

Thus rational points, rational lines, and rational homogeneous conics
can all be treated with the signed-pair arithmetic above. The later
real-name schema explains a separate interpretation for coordinates
or coefficients given by effective real approximations.

